\documentclass[11pt,a4paper]{article}

\usepackage{main-sss}
\usepackage{corrige-sss}

\renewcommand{\matiere}{Physique-Chimie}
\renewcommand{\examen}{DNB}
\renewcommand{\annee}{2026}
\renewcommand{\localisation}{Métropole}
\renewcommand{\titresujet}{Les flotteurs profileurs}
\renewcommand{\niveau}{Brevet}
\renewcommand{\typedocument}{Corrigé}
\renewcommand{\identifiantsujet}{26GENSCMEAG1}

\begin{document}

\pagination
\titre

\q[1 point]
Les flotteurs peuvent mesurer par exemple la pression et la température.

\q[2 points]
L'ion chlorure Cl\textsuperscript{-} a gagné un électron lors de l'ionisation par rapport à l'atome de chlore électriquement neutre.
Il y a donc un électron de plus que de protons, la composition correcte est la 3.

\q[2 points]
On peut déduire des observations des expériences que ce sont les ions qui permettent le passage du courant dans les solutions.

En effet, pour les expériences A et D, dont les solutions ne contiennent que des molécules, le courant ne passe pas et la lampe ne brille pas.

Pour les expériences B et C, les ions chlorure et sodium présents dans la solution permettent le passage du courant et donc la lampe brille.

\q[1 point]
L'appareil qui permet de mesurer l'intensité du courant électrique en ampères est l'ampèremètre.

\q[1 point]
Par lecture graphique, on détermine que pour une intensité du courant électrique mesurée de \si{0,5}{A}, il faut une masse de \si{60}{g} de sel dissous dans un litre d'eau.

\q[2 points]

Le moment entre deux transferts de données correspond à la durée $t$ nécessaire au satellite pour effectuer le tour complet de la Terre.

La distance parcourue est $d = \SI{45000}{km}$ et la vitesse vaut $v = \SI{7,5}{km/s}$.

$t = \frac{d}{v} = \frac{45 000}{7,5} = \SI{6000}{s}$

On montre bien que le transfert de données ne peut s'effectuer que toutes les \si{6000}{s}.

\q[1 point]

Ajouter des satellites permet d'augmenter le nombre de transferts de données, et en les espaçant à des intervalles réguliers, on peut obtenir des mesures beaucoup plus fréquentes.


\end{document}